%\documentclass[12pt,oneside]{amsart}
\documentclass{scrartcl}

%\documentclass{article}

\usepackage{amsthm,amsmath,amssymb}
\usepackage[numbers]{natbib}

%\usepackage[colorlinks]{hyperref}
\usepackage{hyperref}
\hypersetup{
    colorlinks,%
    citecolor=black,%
    filecolor=black,%
    linkcolor=black,%
    urlcolor=black
}
\usepackage{hypernat}
\usepackage[top=2.5cm, bottom=2.5cm, left=2.8cm,
right=2.8cm]{geometry}
\usepackage{graphicx}


%\setlength{\parskip}{0pt}
%\setlength{\parsep}{0pt}
%\setlength{\headsep}{0pt}
%\setlength{\topskip}{0pt}
%\setlength{\topmargin}{0pt}
%\setlength{\topsep}{0pt}
%\setlength{\partopsep}{0pt}

\linespread{1}

%\usepackage{marvosym}

%\textwidth = 6.5 in \textheight = 9.2 in \oddsidemargin = 0.0 in
%\evensidemargin = 0.0 in \topmargin = -0.0 in \headheight = 0.0 in
%\headsep = 0.0in \parskip = 0.0in \parindent = 0.0in
%\def\baselinestretch{0.871}


\newtheorem{theorem}{Theorem}[section]
\newtheorem{proposition}[theorem]{Proposition}
\newtheorem{problem}[theorem]{Problem}
\newtheorem{fact}[theorem]{Fact}
\newtheorem{lemma}[theorem]{Lemma}
\newtheorem{defn}[theorem]{Definition}
\newtheorem{corollary}[theorem]{Corollary}
\newtheorem{remark}{Remark}[section]
\newtheorem{example}[theorem]{Example}

\newcommand{\E}{{\mathbb E}}
\newcommand{\se}{{\mathcal{E}}}
\newcommand{\R}{{\mathbb R}}
\renewcommand{\P}{{\mathbb P}}
\newcommand{\ac}{{\mathcal{A}}}
\newcommand{\A}{{\mathcal{A}}}
\newcommand{\B}{{\mathcal{B}}}
\newcommand{\C}{{\mathcal{C}}}
\newcommand{\M}{{\mathcal{M}}}
\newcommand{\Mf}{{\mathfrak{M}}}
\newcommand{\T}{{\mathcal{T}}}
\newcommand{\Z}{{\mathcal{Z}}}
\newcommand{\K}{{\mathcal{K}}}
\newcommand{\lc}{{\mathcal{L}}}
\newcommand{\Ps}{{\mathcal{P}}}
\newcommand{\G}{{\mathcal{G}}}
\newcommand{\pa}{{\mathring{p}}}
\newcommand{\F}{{\cal F}}
\newcommand{\samp}{{\mathcal{X}}}
\newcommand{\X}{{\mathcal{X}}}
\newcommand{\hi}{{\hat{\Phi}}}
\newcommand{\data}{{\text{data}}}

\newcounter{rcnt}[section]
\renewcommand{\thercnt}{(\roman{rcnt})}

\def\qt#1{\qquad\text{#1}}

\def\argmin{\mathop{\rm argmin}}
\def\argmax{\mathop{\rm argmax}}


\setlength{\parskip}{1.4 \medskipamount}
%\sloppy
%\linespread{1.3}

\begin{document}

\title{STAT 238 - Bayesian Statistics  \\
  Lecture Seventeen} 
\subtitle{Spring 2026, UC Berkeley}

\author{Aditya Guntuboyina}


\date{04 March 2026}

\maketitle

\section{Logistic Regression}

We are again in the
usual regression setting where we observe  data $(y_i,
x_{i1},x_{i2}, \dots, x_{im})$ for $i = 1, \dots, n$. There are $m$
explanatory variables $x_1, \dots, x_m$ and one response variable.
$x_{ij}$ denotes the value of the $j^{th}$ explanatory variable for
the $i^{th}$ individual and $y_i$ is the value of the response
variable for the $i^{th}$ individual. Suppose now that the response
variable is binary i.e., $y_1, \dots, y_n$ take values in $\{0,
1\}$. In this case, the logistic regression model assumes that: 
  \begin{equation*}
    y_i  \overset{\text{independent}}{\sim} \text{Bernoulli} \left(\frac{\exp
        \left(\beta_0 + \beta_1 x_{i1} + \dots + \beta_m x_{im} \right)}{1 + \exp \left(\beta_0 + \beta_1 x_{i1} + \dots + \beta_m x_{im} \right)} 
    \right) ~~\text{for $i = 1, \dots, n$}. 
  \end{equation*}
Letting $x_i = (1, x_{i1}, \dots, x_{im})^T$  and $\beta = (\beta_0,
\beta_1, \dots, \beta_m)^T$, we can write the model also as
  \begin{equation*}
    y_i  \overset{\text{independent}}{\sim} \text{Bernoulli} \left(\frac{\exp
        \left(x_i^T \beta\right)}{1 + \exp \left(x_i^T \beta \right)} 
    \right) ~~\text{for $i = 1, \dots, n$}
  \end{equation*}
  where $\beta$ is the $(m+1) \times 1$ vector with components
  $\beta_0, \beta_1, \dots, \beta_m$.

  This gives the likelihood:
  \begin{align*}
  \text{Likelihood} =\prod_{i=1}^n \left(\frac{\exp
        \left(x_i^T\beta \right)}{1 + \exp \left(x_i^T\beta \right)}
    \right)^{y_i} \left(1 - \frac{\exp
        \left(x_i^T\beta \right)}{1 + \exp \left(x_i^T\beta \right)}
    \right)^{1-y_i} 
  \end{align*}
Suppose $x_1^T, \dots, x_n^T$ form the rows of the $n \times p$
  design matrix $X$ (where $p = m+1$). The unknown parameters in the
  logistic regression model are $\beta_0, \dots, \beta_m$ (note that,
  in contrast to the linear regression model, there is no $\sigma$
  parameter in logistic regression). As in linear regression, we use the prior: 
\begin{equation}\label{unilog}
\beta_0,  \beta_1, \dots, \beta_m \overset{\text{i.i.d}}{\sim} \text{Unif} (-C, C)
\end{equation}
  for a large $C$. The posterior of $\beta$ is then
  \begin{align*}
    f_{\beta \mid \text{data}}(\beta) &\propto \text{Likelihood} \times \text{prior} \\
    &\propto \left[\prod_{i=1}^n \left(\frac{\exp
        \left(x_i^T\beta \right)}{1 + \exp \left(x_i^T\beta \right)}
    \right)^{y_i} \left(1 - \frac{\exp
        \left(x_i^T\beta \right)}{1 + \exp \left(x_i^T\beta \right)}
    \right)^{1-y_i} \right] I\{\beta_0, \beta_1, \dots, \beta_p \in (-C, C)\}
    \\ &= \left[\prod_{i=1}^n \frac{\exp(y_ix_i^T\beta)
    }{1 + \exp(x_i^T\beta)} \right]I\{\beta_0, \beta_1, \dots, \beta_p
         \in (-C, C)\} \\
    &= \left[\exp \left(\ell(\beta) \right) \right]I\{\beta_0, \beta_1, \dots, \beta_p
         \in (-C, C)\}
  \end{align*}
  where
\begin{equation*}
  \ell(\beta) := \sum_{i=1}^n \left[ y_i (x_i^T\beta) - \log \left(1 +
    \exp(x_i^T \beta) \right) \right]. 
\end{equation*}
Note that $\ell(\beta)$ is simply the log-likelihood in this
problem. This posterior density is not in standard form involving some
well-known densities. If $p = 1$ or $p = 2$, 
then this can be plotted. One can use various MCMC techniques to
obtain samples from this posterior.

Instead of MCMC, we shall derive a closed form multivariate normal
approximation that works quite well in practice. This method is called
Laplace Approximation or Posterior Normal Approximation. It turns out
that Bayesian inference from this normal approximation to the
posterior coincides with usual frequentist inference for logistic
regression.

To get the normal approximation, first let us drop the
indicator which will be irrelevant when $C$ is large to get
\begin{align*}
  f_{\beta | Y = y}(\beta) \propto \exp(\ell(\beta)). 
\end{align*}
The normal approximation will be obtained by a second-order Taylor
expansion of $\ell(\beta)$. We shall do the Taylor expansion around
the MLE $\hat{\beta}$ because the posterior is peaked at $\hat{\beta}$
and the high regions of the posterior are most likely very close to
$\hat{\beta}$. Recall that the MLE $\hat{\beta}$ is defined as the
maximizer of the likelihood (or log-likelihood):
\begin{align*}
  \hat{\beta} := \underset{\beta \in \R^p}{\text{argmax}} ~
  \ell(\beta). 
\end{align*}
It is obtained by taking the gradient of the log-likelihood and
solving the equation obtained by setting the gradient to zero. It is
easy to check that the gradient of the log-likelihood is
\begin{equation}\label{grad1}
  \nabla \ell(\beta) = \sum_{i=1}^n \left(y_i - \frac{\exp
        \left(x_i'\beta \right)}{1 + \exp \left(x_i'\beta \right)}
    \right) x_i. 
  \end{equation}
  To get the maximum likelihood estimator $\hat{\theta}$, we need to
  set the gradient above to zero and solve the resulting equation for
  $\theta$. This cannot be done in closed form and the usual method is
  to use an iterative scheme such as Newton's algorithm. The answer
  can be obtained from inbuilt functions in R or Python. More details
  behind the Newton algorithm will be provided a bit later. 

   Coming back to the posterior $\exp(\ell(\beta))$, Taylor expansion of
   $\ell(\beta)$ around the MLE $\hat{\beta}$ gives
   \begin{align*}
     f_{\beta | Y = y}(\beta) \propto \exp(\ell(\beta)) \approx \exp
     \left(\ell(\hat{\beta}) + \left<\nabla \ell(\hat{\beta}), \beta -
     \hat{\beta}\right> + \frac{1}{2} (\beta - \hat{\beta})^T
     H\ell(\hat{\beta}) (\beta - \hat{\beta}) \right)
   \end{align*}
   where $H\ell(\beta)$ denotes the Hessian of $\ell(\beta)$: 
  \begin{equation*}
    H \ell(\beta) = - \sum_{i=1}^n \frac{\exp(x_i' \beta)}{\left(1 +
      \exp(x_i' \beta)\right)^2} x_i x_i'. 
\end{equation*}
 Because the $\ell(\hat \beta)$ term is a constant, it can be ignored
 in proportionality. Also $\nabla \ell(\hat{\beta})$ equals zero. We
 thus have
 \begin{align*}
   f_{\beta \mid Y = y} (\beta) \propto \exp
     \left(\frac{1}{2} (\beta - \hat{\beta})^T
     H\ell(\hat{\beta}) (\beta - \hat{\beta}) \right) = \exp
     \left(-\frac{1}{2} (\beta - \hat{\beta})^T
     \left(-H\ell(\hat{\beta}) \right) (\beta - \hat{\beta}) \right) 
 \end{align*}
We have switched above to $-H\ell(\hat{\beta})$ because this matrix is
positive semi-definite as $\hat{\beta}$ maximizes $\ell(\beta)$. The
above term is simply the unnormalized density of the multivariate
normal distribution with mean $\hat{\beta}$ and covariance matrix $(-H
\ell(\hat{\beta}))^{-1}$. Observe that 
\begin{equation*}
 - H \ell(\hat \beta) =  \sum_{i=1}^n \frac{\exp(x_i' \hat{\beta})}{\left(1 +
      \exp(x_i' \hat{\beta})\right)^2} x_i x_i'. 
\end{equation*}
Now let $W$ denote the $n \times n$ diagonal matrix whose
$i^{th}$ diagonal entry is
\begin{equation*}
  \frac{\exp(x_i' \hat{\beta})}{\left(1 +
      \exp(x_i' \hat{\beta})\right)^2}
\end{equation*}
and also recall again that $X$ is the $n \times p$ matrix with rows
$x_1', \dots, x_n'$. It is then easy to check that
\begin{equation}\label{hess1}
-  H \ell(\hat \beta) = X' W X. 
\end{equation}
The posterior normal approximation is thus
\begin{equation}\label{blogna}
  N(\hat{\beta}, (X' W X)^{-1}). 
\end{equation}
The standard errors corresponding to $\beta_0, \dots, \beta_m$ can
then be obtained by the square roots of the diagonal entries of $(X' W
X)^{-1}$. 

It turns out that Bayesian inference done with the above normal
posterior approximation \eqref{blogna} coincides with the frequentist
inference in 
the logistic regression model. It is easy to check this, say, in
Python using statsmodels
(construct a 95\% credible interval for, say, one of the components of
$\beta$ and then compare it with the frequentist interval). Thus the
usual frequentist inference for the logistic regression model can be
viewed from a Bayesian perspective.

Note that the above analysis
relies on two assumptions: (a) the prior for $\beta$ is assumed to be
uniform on the large cube $(-C, C)^p$, and (b) the posterior is
approximated by a normal distribution. These assumptions may be of
course not reasonable in a particular application. In such a
situation, it is conceptually very clear as to how one would proceed: 
if the normal approximation to the posterior is not accurate, one
needs to work with the actual posterior. If the uniform prior is not
reasonable, one can do the full posterior analysis (or by taking a
normal approximation to the posterior) for a more appropriate prior. 

\subsection{Details behind the Newton Algorithm for computing the MLE}

The MLE $\hat{\beta}$ of $\beta$ is the maximizer of $\ell(\beta)$. The maximizer of
$\ell(\beta)$ cannot be computed in closed form. Newton's
method is commonly used for maximizing $\ell(\beta)$. Newton's method
uses the iterative scheme
\begin{equation}\label{ish}
  \beta^{(m+1)} = \beta^{(m)} - \left(H\ell(\beta^{(m)})\right)^{-1}
  \nabla \ell(\beta^{(m)}) 
\end{equation}
As we saw in \eqref{grad1},
\begin{equation*}
  \nabla \ell(\beta) = \sum_{i=1}^n \left(y_i - \pi_i \right) x_i
\end{equation*}
where $\pi_i$ is given by
\begin{equation*}
  \pi_i = \frac{\exp(x_i' \beta)}{1 + \exp(x_i' \beta)}
\end{equation*}
Letting $\pi$ be the $n \times 1$ vector with entries $\pi_1, \dots,
\pi_n$, we can write $\nabla \ell(\beta)$ in matrix notation as (note
that $X$ has rows $x_1^T, \dots, x_n^T$ or, equivalently, $X^T$ has
columns $x_1, \dots, x_n$):
\begin{equation*}
\nabla  \ell(\beta) = X^T \left(y - \pi \right)
\end{equation*}
where, as usual in regression, $y$ denotes the $n \times 1$  vector of
response values. Also from \eqref{hess1}, we can write
\begin{equation*}
  H\ell(\beta) = - X^T W X
\end{equation*}
where $W$ is the $n \times n$ diagonal matrix whose $i^{th}$ diagonal
entry is $\pi_i(1 - \pi_i)$. Newton's iterative scheme~\eqref{ish}
therefore becomes 
\begin{equation*}
  \beta^{(m+1)} = \beta^{(m)} + (X^T W X)^{-1} X^T (y - \pi). 
\end{equation*}
This can be rewritten as 
\begin{equation}\label{wlse}
  \beta^{(m+1)} = (X^T W X)^{-1} X^T W Z
\end{equation}
where 
\begin{equation}\label{zdef}
  Z = X \beta^{(m)} + W^{-1}(y - \pi). 
\end{equation}
The method of obtaining the MLE $\hat{\beta}$ therefore proceeds iteratively as
follows. First have an initial estimate of $\beta$. Call this initial
estimator $\hat{\beta}^{(0)}$. Use this 
estimator to calculate $\pi_i$ via 
\begin{equation*}
\pi_i = \frac{\exp(\hat{\beta}^{(0)}_0 + \hat{\beta}^{(0)}_1 x_{i1} + \dots +
  \hat{\beta}^{(0)}_p x_{ip})}{1 + \exp(\hat{\beta}^{(0)}_0 +
  \hat{\beta}^{(0)}_1 x_{i1} \dots + \hat{\beta}^{(0)}_p x_{ip})}.        
\end{equation*}
Use these values of $\pi_i$ to create the response variable values $Z_i$
via~\eqref{zdef} and also use values of $\pi_i$ to construct the matrix
$W$. With $Z$ and $W$, we can estimate $\beta$ via 
\begin{equation*}
  \hat{\beta}^{(1)} = (X^T W X)^{-1} X^T W Z. 
\end{equation*}
Now replace the initial estimator $\hat{\beta}^{(0)}$ by
$\hat{\beta}^{(1)}$ and repeat this process. Keep repeating this until
two successive estimates $\hat{\beta}^{(m)}$ and $\hat{\beta}^{(m+1)}$
do not change much. At that point, stop and report the estimate of
$\beta$ in the logistic regression model as $\hat{\beta}^{(m)}$. 

The expression $(X^T W X)^{-1} X^T W Z$ is reminiscent of the usual
$(X^T X)^{-1} X^T y$ which is the usual estimate of $\beta$ in the
linear model. In fact, this is the least squares estimate in a
weighted least squares model. 




%\begin{thebibliography}{99}

%\bibitem[MacKay and Peto(1995)]{MacKay1995HierarchicalDirichlet}
%D.~J.~C. MacKay and L.~C. Peto,
%\newblock ``A hierarchical Dirichlet language model,''
%\newblock \emph{Natural Language Engineering}, vol.~1,
%no.~3, pp.~289--308, 1995.

%\bibitem[Rubin(1981)]{Rubin1981BayesianBootstrap}
%D.~B. Rubin,
%\newblock ``The Bayesian bootstrap,''
%\newblock \emph{The Annals of Statistics}, vol.~9,
%no.~1, pp.~130--134, 1981.
  
%\bibitem[Fisher(1934)]{Fisher1934}
%R.~A. Fisher,
%\newblock ``Probability, likelihood and quantity of information in the logic of
%uncertain inference,''
%\newblock \emph{Proceedings of the Royal Society of London. Series~A,
%Containing Papers of a Mathematical and Physical Character}, vol.~146,
%no.~856, pp.~1--8, 1934.

%\bibitem[Efron(2010)]{Efron}
%Efron, B. (2010)
%\textit{Large-scale inference: empirical Bayes methods for estimation,
%testing, and prediction}. Cambridge University Press, 2010.

%\bibitem[Shumway and Stoffer(2017)]{ShumwayStoffer}
%Shumway, R. H. and Stoffer, D. S. (2017)
%\textit{Time Series Analysis and Its Applications: With R
%  Examples}. Springer, 4th edition. 
  
%     \bibitem[Jaynes(2003)]{Jaynes} Jaynes, E. T, (2003)
%  \textit{Probability theory: the logic of science}. Cambridge
%  University Press, 2003.
 
%  \bibitem[Sinai(1992)]{Sinai} Sinai, Y. G, (1992)
%  \textit{Probability theory: an introductory course}. Springer
%  Verlag, 1992.  

%\bibitem[{deGroot(1986)}]{DeGrootBlackwell}DeGroot, M. H. (1986)
%       A conversation with David Blackwell.
%\newblock \emph{Statistical Science}, \textbf{1}, 40--53.

%         \bibitem[Mosteller(1987)]{Mosteller} Mosteller, F, (1987)
%  \textit{Fifty challenging problems in probability with
%    solutions}. Courier Corporation, 1987.

%    \bibitem[MacKay(2003)]{MacKay} MacKay, D. J. C., (2003)
%  \textit{Information theory, inference, and learning
%  algorithms}. Cambridge University Press, 2003.

%\bibitem[{Lindley(2013)}]{Lindley}Lindley, D. V. (2013)
%   \textit{Understanding Uncertainty}. John Wiley and sons, 2013.


%\end{thebibliography}







\end{document}

%%% Local Variables:
%%% mode: latex
%%% TeX-master: t
%%% End:
